\section{Equilibrium Analysis}
\label{sec:analysis}

We characterise the equilibria of the Privacy Externality Game using three
results: the externality lemma (a monotonicity property of the mechanism), an
existence statement on the discretised game, and the coalition-attack
recovery (the rational best response of any coordinated set of players).

\subsection{Monotonicity of the externality}

\begin{lemma}[Privacy externality]
\label{lem:externality}
Consider the sampling-based iDP mechanism of \cref{def:sampling-idp} on a
feasible profile $\boldsymbol{\eps}$ at which constraints (C1) bind for every
group and (C2) binds at $E_b$. Then for every pair $i \neq j$,
\begin{equation}
\frac{\partial A_i(\boldsymbol{\eps})}{\partial \eps_j} \;\leq\; 0,
\label{eq:externality-lemma}
\end{equation}
with equality iff group $i$'s sampling rate $p_i$ does not respond to a
change in $\eps_j$.
\end{lemma}

\begin{proof}
Work in the R\'enyi-DP form of (C1) of \citet{mironov2017renyi}: at order
$\alpha$, the RDP of an $I$-fold $(p_g,\sigma)$-subsampled Gaussian is at
most $\rho_g = 2I\alpha\, p_g^2/\sigma^2$, with equality up to higher-order
terms when $p_g$ is small (Mironov, Theorem~5). Setting this equal to the
target RDP $\rho_g(\eps_g, \delta)$ gives
\[
p_g \;=\; \frac{\sigma}{\sqrt{2I\alpha}}\,\sqrt{\rho_g(\eps_g, \delta)}.
\]
Substituting into (C2) yields the implicit equation for $\sigma$:
\begin{equation}
\sum_{g}\, \frac{|G_g|}{M}\,\frac{\sigma}{\sqrt{2I\alpha}}\,
\sqrt{\rho_g(\eps_g,\delta)} \;=\; \frac{E_b}{M}.
\label{eq:sigma-implicit}
\end{equation}
Since $\rho_g(\cdot, \delta)$ is monotone-decreasing in $\eps_g$ (smaller
budget = smaller RDP), the left-hand side is decreasing in $\sigma$ at fixed
$\eps_g$ if and only if it is increasing in any $\eps_j$. Implicit
differentiation gives $\partial \sigma / \partial \eps_j > 0$ when (C2)
binds. Recall that $p_i$ is increasing in $\sigma$ at fixed $\eps_i$.
Consequently $\partial p_i / \partial \eps_j > 0$ whenever group $i$ is
non-degenerate.

The MIA advantage of an $I$-fold $(p_i,\sigma)$-subsampled Gaussian is
non-decreasing in $p_i$ (more sampling exposes more data) and non-increasing
in $\sigma$ (more noise hides more), as is the privacy profile of the
subsampled Gaussian \cite{balle2020privacy, dong2022gaussian}. Lowering
$\eps_j$ reverses the signs: $\partial \sigma / \partial \eps_j < 0$ and
$\partial p_i / \partial \eps_j < 0$. Net effect on $A_i$ along the feasible
surface is non-positive, which is~\eqref{eq:externality-lemma}.
\end{proof}

\Cref{lem:externality} formalises the empirical observation of
\citet{kaiser2026your}: when any player tightens their iDP budget, the
mechanism's noise multiplier $\sigma$ shifts to keep (C2) binding, and every
other player's realised vulnerability $A_i$ rises. The externality is
\emph{monotone} in every other player's budget.

In the decomposition of \cref{sec:decomposition}, the lemma equivalently
states that \emph{the slack $\bar c(\eps_i) - A_i(\boldsymbol{\eps})$
shrinks in every other player's $\eps_j$}. In words, the
$(\eps_i,\delta)$-DP worst-case ceiling $\bar c(\eps_i)$ is unchanged
(it depends only on $\eps_i$), but the realised advantage $A_i$ climbs
toward that ceiling as others tighten their budgets. This is the
quantity that the proposed $(\eps_i, \delta, \Delta)$-iDP
contract~\cite{kaiser2026your, kaissis2024beyond} would cap, and our
Lemma motivates such a cap formally: without it, $\mathrm{slack}_i$ has no
floor and the externality can drive any player's realised advantage
arbitrarily close to the worst case.

\paragraph{Minority and majority intuition.}
The externality has an asymmetric \emph{group-size} structure. Because (C2)
enforces a fixed expected batch size, a player with small $|G_i|$ but
small $\eps_i$ contributes little to the batch and forces the
larger-$|G_j|$ players' sampling rates upward. The minority therefore
\emph{exports} its privacy demand onto the majority. This is the formal
content of the empirical observation in \cite{kaiser2026your}
that ``the proportion of low-budget players matters as much as their
$\eps$''.

\subsection{Existence and structure of pure Nash equilibria}

\begin{proposition}[NE existence on the discretised game]
\label{prop:existence}
The discretised PEG on the action grid $\mathcal{A}_K =
\{\eps^{\min},\dots,\eps^{\max}\}\subset[\eps^{\min},\eps^{\max}]$
with utilities~\eqref{eq:utility} is a finite normal-form game restricted to
the (closed) feasible set. By the Nash existence
theorem~\cite{nash1950equilibrium}, every such game admits at least one
mixed-strategy NE. Pure NE may or may not exist, depending on
$(B_i, C_i, |G_i|, V, K)$.
\end{proposition}

We do not claim existence on the continuous action space without further
quasi-concavity assumptions; in particular, our utility~\eqref{eq:utility}
is not generally concave in $\eps_i$ because $A_i$ is the trade-off
function of a subsampled Gaussian, whose dependence on $\eps_i$ through
$(\sigma, p_i)$ is non-trivial. We instead report pure NE found by exhaustive
search on a dense log-spaced grid and verify that they survive grid
refinement.

\paragraph{What the externality lemma implies for NE structure.}
Lemma~\ref{lem:externality} says lowering another player's budget weakly
raises one's own realised vulnerability. If players' utilities place
non-trivial weight on $A_i$ ($C_i > 0$), each player has an incentive to
\emph{shift the externality onto others} by lowering their own budget---an
incentive that is intrinsic to the mechanism. The opposing incentive comes
from $V(\sigma)$: lowering own $\eps$ pushes $\sigma$ upward, lowering the
model utility for everyone (including oneself). The equilibrium balances
these two incentives.

We observe in our experiments (\cref{sec:experiments}) that the resulting
balance is typically \emph{asymmetric}: in the symmetric 2-player setting,
pure NEs sit off the diagonal of the action grid. The intuition is the
externality lemma's converse: at any symmetric profile $(\eps,\eps)$, both
players face an identical incentive to deviate; since both deviations cannot
be unilaterally absorbed, the equilibrium is at an asymmetric profile where
one player carries a higher $A$ in exchange for a higher $V$.

\subsection{Coalition Stackelberg and the attack recovery}

\begin{proposition}[Coalition best response recovers the attack]
\label{prop:coalition-attack}
Fix a target $t$ and a coalition $C$ of size $k$, with $\lambda \to \infty$
in Definition~\ref{def:coalition} (the coalition cares only about the
target's vulnerability). The coalition's best response is $\eps_i =
\eps^{\star}_k$ for every $i \in C$, where $\eps^{\star}_k$ is the smallest
admissible budget that keeps the coalition's mechanism solution feasible
(i.e., $p_i > 0$ for every $i \in C$). Furthermore, the target's MIA
advantage $A_t$ is non-decreasing in $k$.
\end{proposition}

\begin{proof}[Proof sketch]
By \cref{lem:externality}, $\partial A_t / \partial \eps_i \leq 0$ for every
$i \in C$ on the feasible region. Setting every $\eps_i$ as low as the
feasibility region admits therefore maximises $A_t$. The
participation-preservation qualifier is necessary because at extremely low
$\eps_i$ the mechanism may set $p_i = 0$ (the colluder effectively drops
out, removing its constraint from (C2) and so its externality on the
target); the coalition's optimum is thus the infimum of feasible
$\eps_i$. Monotonicity in $k$ is the same lemma applied to adding a member
to $C$: doing so strictly enlarges the set of budgets that pull $\sigma$
downward, raising $A_t$.
\end{proof}

Proposition~\ref{prop:coalition-attack} establishes that the budget-
manipulation attack of \citet{kaiser2026your} is precisely the
\emph{rational} best response of any coordinated subset of data holders,
regardless of malicious intent. Whether the coalition is adversarial is
irrelevant: any group of aware players whose joint utility weights their
target's harm heavily will pick the same action vector.

\subsection{An aside on awareness}

\Cref{lem:externality} requires that players compute (or approximate) the
realised $A_i$. A natural question is: what about a player whose only
information is the formal $(\eps_i, \delta)$-iDP contract? Such a player
treats $\eps_i$ as a direct privacy proxy and plays the budget that
maximises $V(\sigma) - g(\eps_i)$ for some monotone $g$. Their best response
does not depend on $\eps_{-i}$ at all (the externality is invisible), so
their action is mechanically the smallest $\eps_i$ that does not break
their own $V$. This is the \emph{rational} action of a naive player, and it
is precisely the configuration in which the attack of
\citet{kaiser2026your} succeeds against them: by Lemma~\ref{lem:externality},
naive players' choices feed the externality on everyone else without their
realising. The contribution of this paper is to show that even
\emph{aware} players---those who can compute $A_i$ explicitly---arrive at
asymmetric equilibria that concentrate vulnerability on a subset of players.
The externality cannot be undone by individual rationality alone.
